\chapter{Platelet-Rich Plasma Injection}

\section*{What Is This Test or Procedure?}
A platelet-rich plasma (PRP) injection is a procedure in which a concentrate of the patient's own platelets is injected into a tendon, ligament, joint, or other tissue to promote healing and reduce pain.
\section*{Alternative Names}
PRP injection; autologous platelet concentrate injection; platelet-rich plasma therapy
\section*{Principle}
Platelets release growth factors and signaling proteins that recruit healing cells and stimulate tissue repair. Concentrating the patient's platelets and delivering them to the injury site amplifies this natural healing response.
\section*{History}
The regenerative potential of platelet concentrates was recognized in the 1970s and popularized in orthopedic medicine in the 2000s as an option for tendinopathy and osteoarthritis.
\section*{Why This Test or Procedure Is Ordered}
Ordered for chronic tendinopathy, such as lateral epicondylitis or Achilles tendinopathy, selected ligament and muscle injuries, and knee osteoarthritis, when symptoms persist despite conservative treatment.
\section*{Is This a Primary or Secondary Test or Procedure}
A secondary therapeutic procedure, used after conservative measures have not fully succeeded and before or in place of more invasive treatments.
\section*{How This Test or Procedure Is Performed}
Blood is drawn from the patient and centrifuged to separate platelets from red and white blood cells. The platelet concentrate is then injected into the affected tissue, often with ultrasound guidance, either as a single treatment or a series.
\section*{What Preparation Is Needed}
Informed consent and review of bleeding risk. Patients are advised to stop anti-inflammatory medications such as NSAIDs before the injection because they may interfere with the healing response.
\section*{Contraindications and Precautions}
Contraindications include local infection, bleeding disorders, and active systemic infection. Precautions include pregnancy, cancer at the treatment site, and the use of medications that affect platelets.
\section*{Risks}
Pain and soreness at the injection site, bruising, and rarely infection or nerve injury. There is limited evidence of benefit for some conditions, so expectations should be realistic.
\section*{What Happens After the Test or Procedure}
The patient rests the treated area briefly and gradually resumes activity. Several weeks may be required before improvement, and more than one injection is sometimes needed.
\section*{Understanding Your Results}
Response is judged by reduction in pain and improvement in function over the weeks after injection. Lack of improvement may indicate the condition will not respond to this treatment.
\section*{Normal Results}
Reduced pain and improved function with no significant adverse reaction at the injection site.
\section*{Follow-up Tests and Procedures}
Patients are reassessed several weeks after the injection; repeat injections may be offered, and persistent symptoms may lead to other treatments such as surgery.
\section*{References}
\begin{enumerate}
  \item MedlinePlus. Platelet-rich plasma injection. U.S. National Library of Medicine; 2025.
  \item Current Medical Diagnosis and Treatment. McGraw-Hill; 2025.
\end{enumerate}

\clearpage
